\chapter{Conclusiones y trabajo futuro}

\section{Conclusiones técnicas}
% \subsection{Aportaciones principales}
% \subsubsection{Impacto en el proyecto}

\section{Cumplimiento de objetivos}
% \subsection{Objetivos alcanzados}
% \subsubsection{Objetivos parcialmente alcanzados}

\section{Trabajo futuro}

% Modulo scraping actualizar GTFS
% \subsection{Integración completa IA-web}
% \subsubsection{Endpoint y visualización de probabilidades}

% \subsection{Mejoras de modelado}
% \subsubsection{Comparativa ampliada con DNN}

% \subsection{Evolución del simulador}
% \subsubsection{Escalado y nuevas funcionalidades}

% Trabajar con la probabilidad permite posteriormente obtener cuotas de mercado que pueden ser muy interesantes para el simulador, porque con el simulador quieres ver si modificas algo en el sistema, como cambian las cutoas del mercado (porcentaje de usuarios de cada medio de transporte). Con el simulador haces una simulación de cambiar los valores de los atributos del usuario (cambios en precio de bus, predecir el comportamietno de como esos aspectos afectan a la cuota de mercado de esos perfils de usuario). Trabajar con las probabilidad en vez de con la elección final puedes trabajar mejor.

% Cerrar capitulo con ejemplos pantallazo, o incluso una sección adicional de mini-demo.